\documentclass[11pt,a4paper]{article}
\usepackage[utf8]{inputenc}
\usepackage[T1]{fontenc}
\usepackage[french]{babel}
\usepackage{geometry}
\usepackage{xcolor}
\usepackage{hyperref}
\usepackage{titlesec}

\geometry{margin=1in}
\definecolor{titleblue}{RGB}{0,51,102}

\hypersetup{
    colorlinks=true,
    linkcolor=titleblue,
    urlcolor=titleblue,
    pdftitle={Lettre de Motivation - Ismail AISSAOUI},
}

\begin{document}

\noindent \textbf{Ismail AISSAOUI} \\
Ingénieur d'État en Intelligence Artificielle (ESI-SBA) \\
Email: ismail.aissaoui.pro@gmail.com \\
Téléphone: +213 660 70 77 96 \\
LinkedIn: https://ismail-aissaoui.vercel.app

\bigskip

\begin{flushright}
    Au Comité de Sélection du Doctorat \\
    \textbf{Programme EDT - CEA List / IRIT} \\
    À l'attention de : Pr. Brahim Hamid et Dr. Marcos Didonet Del Fabro \\
    \medskip
    1er mai 2026
\end{flushright}

\bigskip

\noindent \textbf{Objet : Candidature au doctorat : Ingénierie de la sécurité assistée par l'IA pour les Jumeaux Numériques (Rang 5 - PC3)}

\bigskip

Madame, Monsieur,

Ingénieur d'État en Intelligence Artificielle diplômé de l'ESI-SBA, je vous adresse avec enthousiasme ma candidature pour le poste de doctorat intitulé « Ingénierie de la sécurité assistée par l'IA pour les Jumeaux Numériques » au sein du programme national Engineering Digital Twin (EDT). Cette candidature est motivée par la forte adéquation de mon profil avec vos recherches sur l'intelligence relationnelle et l'ingénierie dirigée par les modèles (MDE).

Actuellement en fin de cycle ingénieur chez Sonatrach, je travaille sur le développement de **Jumeaux Numériques informés par la physique** pour la surveillance d'infrastructures critiques. Au-delà des aspects purement physiques, j'ai acquis une solide base en **Réseaux de Neurones sur Graphes (GNN)** que j'ai appliquée à la détection d'anomalies structurelles et à l'extraction de connaissances. Je suis particulièrement motivé par l'idée d'appliquer ces compétences aux défis de l'injection automatique de sécurité au sein des frameworks MDE, garantissant ainsi des jumeaux numériques « sécurisés par conception » (secure-by-design).

Votre projet au CEA List et à l'IRIT représente une opportunité unique de travailler à l'intersection de la cybersécurité et de l'ingénierie système. Je suis convaincu que ma maîtrise de PyTorch et ma compréhension des problématiques industrielles réelles me permettront de contribuer de manière significative à vos objectifs de synchronisation bidirectionnelle (round-trip engineering) pour la sécurité.

Je suis un chercheur autonome et rigoureux, capable de traduire des besoins complexes en solutions algorithmiques robustes pour la **plateforme Artemis**.

Je vous remercie de l'attention que vous porterez à ma candidature et reste à votre entière disposition pour un entretien.

Dans l'attente de votre réponse, je vous prie d'agréer, Madame, Monsieur, l'expression de mes salutations distinguées.

\bigskip

\noindent Ismail AISSAOUI

\end{document}
